\section{Bayesian models for the dark figure}\label{sec:models}

\Cref{sec:identification} isolates what any estimator of the dark figure must assume. This section turns that analysis into estimation. We develop four Bayesian components. Each addresses one way in which the single-domain model of \cref{sec:framework} is too simple for antisemitic hate crime. Incidents differ in how reportable they are (\cref{sec:models-lcmcr}). Some of that variation is visible in recorded incident characteristics (\cref{sec:models-cov}). Most state--years contain few incidents, and the coverage of the police list changes abruptly over time (\cref{sec:models-sae}). Finally, the capture histories are not observed: they are estimated by record linkage (\cref{sec:models-link}). In every component the unobserved count $\nzero$ is a latent variable that is sampled together with the parameters, so every source of uncertainty reaches the reported interval. None of the components is new on its own. Our contribution is to assemble them for this application and to make explicit where each one obtains its identifying information. Posterior computation is described in \cref{app:computation}.

\subsection{Latent-class MSE with a Dirichlet-process prior}\label{sec:models-lcmcr}

Consider two incidents in the same city in the same month. In the first, a synagogue is attacked and a congregant is assaulted. The police respond, the local press reports the attack, the NGO audit records it from the press and from the congregation, and the community security organisation, which works with the congregation, records it as well. In the second, a man wearing a kippah is sworn at on a street corner. He tells his family and perhaps files a report on a community portal. No other list hears of it. Incidents of the first kind are recorded by several lists; incidents of the second kind are recorded by at most one. Within each kind, whether one list records an incident may be nearly unrelated to whether another does. In a population that mixes the two kinds, however, the lists are positively dependent. A police record is evidence that the incident is of the visible kind, and therefore that the NGO also recorded it. \Cref{prop:id-hetero} formalises this mechanism. With two lists, such dependence biases the Lincoln--Petersen estimator downward (\cref{prop:direction}). With three or more lists, the sign of the bias of the model without the highest-order interaction is not determined by heterogeneity alone (\cref{prop:id-hetero}), which is one reason we model the heterogeneity directly.

Latent-class models represent such heterogeneity by a discrete \emph{reportability type} $z_i\in\{1,\dots,J\}$, within which the lists record independently. \Citet{manriquevallier2016} placed a Dirichlet-process prior on the distribution of types, so that the number of types need not be fixed in advance. We use his model (often called LCMCR) with the truncated stick-breaking representation of the Dirichlet process \citep{sethuraman1994,ishwaran2001}, truncated at $J^\star$ classes:
\begin{equation}\label{eq:lcmcr}
\begin{aligned}
  Y_{ik}\mid z_i=j &\overset{\text{ind}}{\sim}\Bern(\lambda_{jk}),\quad k\in\Kset,
  &\qquad
  z_i &\overset{\text{iid}}{\sim}\operatorname{Categorical}(\pi_1,\dots,\pi_{J^\star}),\quad i=1,\dots,N,\\
  \pi_j &= V_j\textstyle\prod_{l<j}(1-V_l),
  &
  V_j\mid\alpha &\overset{\text{iid}}{\sim}\operatorname{Beta}(1,\alpha)\ \ (j<J^\star),\quad V_{J^\star}=1,\\
  \lambda_{jk} &\overset{\text{iid}}{\sim}\operatorname{Beta}(a_\lambda,b_\lambda),
  &
  \alpha &\sim\operatorname{Gamma}(a_\alpha,b_\alpha),\qquad p(N)\propto 1/N .
\end{aligned}
\end{equation}
We use $a_\lambda=b_\lambda=1$ and $a_\alpha=b_\alpha=0.25$ (shape and rate), a diffuse choice whose influence we examine in \cref{app:computation}. Integrating out $z_i$, the histories are independent and identically distributed, as in model (M) of \cref{def:fw-models}, with cell probabilities
\begin{equation}\label{eq:models-cell}
  \pi_{\h}(\psi)=\sum_{j=1}^{J^\star}\pi_j\prod_{k=1}^{K}\lambda_{jk}^{h_k}(1-\lambda_{jk})^{1-h_k},
  \qquad \h\in\Hset,\qquad \psi=(\pi_1,\dots,\pi_{J^\star},\bm\lambda).
\end{equation}
Bold subscripts index histories and plain subscripts index classes. This is the mixture representation of \cref{sec:fw-sampling} with a discrete mixing distribution. In particular, $\pi_{\zero}(\psi)=\sum_j\pi_j\prod_k(1-\lambda_{jk})$ is the probability that an incident is missed by every list. The visible type of the example above corresponds to a class $j$ whose $\lambda_{j\cdot}$ are all large. The low-visibility type corresponds to a class with small $\lambda_{j\cdot}$, and such a class contributes most of $\pi_{\zero}$.

The prior $p(N)\propto1/N$ has a convenient property, which the next lemma records together with a one-parameter family of alternatives that we use for sensitivity analysis. The calculation is elementary; we state it because it underlies both the sampler and the Bayesian reading of the Horvitz--Thompson estimator in \cref{sec:models-cov}.

\begin{lemma}[The unobserved count]\label{lem:models-negbin}
Let $(n_{\h})_{\h\in\Hset}\mid N,\psi\sim\Mult(N,\bm\pi(\psi))$ for any parametric capture model with $\pi_{\zero}(\psi)<1$. Suppose that only $\bm n_+$ is observed, with $n\ge1$. Let $N$ and $\psi$ be independent a priori, with $p_\kappa(N)\propto\Gamma(N+\kappa)/\Gamma(N+1)$ on $N\ge1$, for some $\kappa\ge0$; $\kappa=0$ gives $p(N)\propto1/N$ and $\kappa=1$ the flat prior. Then:
\begin{enumerate}[(i)]
  \item $\nzero\mid\psi,\bm n_+\sim\operatorname{NegBin}\big(n+\kappa,\,1-\pi_{\zero}(\psi)\big)$, the number of failures before the $(n+\kappa)$-th success when the success probability is $1-\pi_{\zero}$. Hence
  $\E[N\mid\psi,\bm n_+]=\{n+\kappa\,\pi_{\zero}(\psi)\}/\{1-\pi_{\zero}(\psi)\}$.
  \item $p(\psi\mid\bm n_+)\propto p(\psi)\,\{1-\pi_{\zero}(\psi)\}^{-\kappa}\prod_{\h\in\Hobs}\pi^c_{\h}(\psi)^{n_{\h}}$, with $\pi^c_{\h}=\pi_{\h}/(1-\pi_{\zero})$ as in \cref{lem:fw-factor}.
\end{enumerate}
In particular, under $p(N)\propto1/N$ the marginal posterior of $\psi$ is the posterior obtained from the conditional likelihood alone, and $\E[N\mid\psi,\bm n_+]=n/\{1-\pi_{\zero}(\psi)\}$.
\end{lemma}

\begin{proof}
By \eqref{eq:fw-likM}, as a function of $\nzero=N-n$ the likelihood is proportional to $N!\,\pi_{\zero}^{\nzero}/\nzero!$. Multiplying by $p_\kappa(N)\propto\Gamma(N+\kappa)/N!$ gives
$p(\nzero,\psi\mid\bm n_+)\propto p(\psi)\,\{\Gamma(n+\nzero+\kappa)/\nzero!\}\,\pi_{\zero}^{\nzero}\prod_{\h\in\Hobs}\pi_{\h}^{n_{\h}}$.
As a function of $\nzero$, this is the kernel of the stated negative binomial, which proves (i); the mean follows from $\E\nzero=(n+\kappa)\pi_{\zero}/(1-\pi_{\zero})$. For (ii), sum over $\nzero\ge0$ using the negative-binomial series $\sum_{x\ge0}\Gamma(r+x)\,q^{x}/x!=\Gamma(r)(1-q)^{-r}$, valid for $r>0$ and $0\le q<1$, with $r=n+\kappa$. This gives $p(\psi\mid\bm n_+)\propto p(\psi)\,(1-\pi_{\zero})^{-(n+\kappa)}\prod_{\h\in\Hobs}\pi_{\h}^{n_{\h}}$. Then substitute $\pi_{\h}=(1-\pi_{\zero})\pi^c_{\h}$ and use $\sum_{\h\in\Hobs}n_{\h}=n$.
\end{proof}

\Cref{lem:models-negbin} is the engine of data augmentation. Given $\psi=(\pi,\bm\lambda)$, the unseen incidents can be generated exactly. First draw $\nzero$ from the negative binomial with success probability $1-\sum_j\pi_j\prod_k(1-\lambda_{jk})$. Then allocate the $\nzero$ unseen incidents to classes with probabilities proportional to $\pi_j\prod_k(1-\lambda_{jk})$. The completed data form a fully observed latent-class sample of size $N=n+\nzero$, and every remaining full conditional is conjugate (\cref{app:computation}). Uncertainty about $N$ therefore enters every sweep of the sampler rather than being added afterwards. Part (ii) shows what the choice of prior on $N$ does. Under $1/N$ the capture parameters are learned exactly as from the conditional likelihood in \eqref{eq:fw-likM}, with no pull toward small or large populations. Under the flat prior ($\kappa=1$), the posterior of $\psi$ is tilted by $(1-\pi_{\zero})^{-1}$ toward low detectability, and hence toward larger $N$. This tilt matters most when $\pi_{\zero}$ is poorly determined, which is exactly the regime of a large dark figure.

\begin{remark}[What the latent classes identify]\label{rem:models-lc-id}
Local independence fixes every interaction \emph{within} a class at zero. The mixture, however, generally has nonzero interactions of all orders, including a nonzero top interaction. Each posterior draw of $\psi$ therefore implies a value of $\omega$, obtained by applying \eqref{eq:fw-omega} to $\bm\pi(\psi)$. We report the posterior of this implied $\omega$. It places the LCMCR estimate on the sensitivity scale of \cref{thm:closed,cor:bounds}, so that a reader can see how far from the no-$K$-way-interaction model the latent-class fit has moved, and in which direction. The observed table does not identify $\omega$ (\cref{thm:nonid}). The posterior of the implied $\omega$ is therefore a consequence of the latent-class structure and the prior on $\psi$, not a measurement. Two results sharpen this caveat. First, different heterogeneity distributions can fit the observed cells equally well and imply very different $N$ \citep{link2003}. Second, generic identifiability of latent-class parameters requires the number of binary indicators to grow with the logarithm of the number of classes \citep{allman2009}; with $K=4$ lists, only a few classes are identified even from the complete table, and here the zero cell is missing. The Dirichlet-process prior concentrates on parsimonious class structures, and that concentration is part of the identifying assumption. For this reason we always report LCMCR next to the closed-form sensitivity analysis and the heterogeneity-robust lower bound of \cref{thm:chao}, and never alone.
\end{remark}

\subsection{Covariate-dependent capture}\label{sec:models-cov}

Some heterogeneity is observed. Lists record characteristics of an incident: the offence type (assault, vandalism, threat, harassment), the target (a person or an institution), the setting (a place of worship, a school or campus, a street, a home) and the urbanicity of the location. Let $\bm x_i$ collect such characteristics, harmonised across lists (\cref{sec:plan}). Consider the model
\begin{equation}\label{eq:models-huggins}
  \logit p_{ik}=\beta_{0k}+\bm x_i^{\top}\bm\beta_k,\qquad
  Y_{i1},\dots,Y_{iK}\ \text{independent given }\bm x_i,
  \qquad p_i(\bm\beta)=1-\prod_{k=1}^{K}(1-p_{ik}),
\end{equation}
where $p_i$ is the probability that incident $i$ is recorded by at least one list. Covariates are observed only for recorded incidents. \Citet{huggins1989} and \citet{alho1990} therefore condition on being recorded. The conditional likelihood and the resulting Horvitz--Thompson estimator are
\begin{equation}\label{eq:models-ht}
  L_c(\bm\beta)=\prod_{i=1}^{n}\frac{\prod_{k}p_{ik}^{y_{ik}}(1-p_{ik})^{1-y_{ik}}}{p_i(\bm\beta)},
  \qquad
  \Nhat_{\mathrm{HT}}=\sum_{i=1}^{n}\frac{1}{p_i(\widehat{\bm\beta})},
\end{equation}
where $\widehat{\bm\beta}$ maximises $L_c$. Given the true $\bm\beta$, the estimator is unbiased, because $\E\sum_{i\in U}\ind{\h_i\ne\zero}/p_i=N$. Huggins' variance estimator adds the estimation error of $\widehat{\bm\beta}$ to the design variance:
\begin{equation}\label{eq:models-hugvar}
  \widehat{\Var}(\Nhat_{\mathrm{HT}})=\sum_{i=1}^{n}\frac{1-\hat p_i}{\hat p_i^{2}}+\widehat{\bm D}^{\top}\widehat{\bm I}^{-1}\widehat{\bm D},
  \qquad
  \widehat{\bm D}=\sum_{i=1}^{n}\frac{\partial p_i^{-1}}{\partial\bm\beta}\bigg|_{\widehat{\bm\beta}},
\end{equation}
where $\widehat{\bm I}$ is the observed information of $\log L_c$.

\paragraph{A Bayesian reading} \Cref{lem:models-negbin} turns \eqref{eq:models-ht} into a posterior statement. Group the incidents by covariate profile, and in the stratum with profile $\bm x$ place the prior $p(N_{\bm x})\propto1/N_{\bm x}$. For continuous covariates, each recorded incident forms its own stratum. Within a stratum the model is a special case of the lemma with $\pi_{\zero}=1-p(\bm x)$. Hence $N_{\bm x}-n_{\bm x}\mid\bm\beta\sim\operatorname{NegBin}(n_{\bm x},p(\bm x))$, which is the law of a sum of $n_{\bm x}$ independent geometric variables. Collecting strata,
\begin{equation}\label{eq:models-geom}
  N\mid\bm\beta,\text{data}\ \overset{d}{=}\ n+\sum_{i=1}^{n}G_i,\quad G_i\overset{\text{ind}}{\sim}\operatorname{Geom}\big(p_i(\bm\beta)\big)\ \text{on }\{0,1,\dots\},
  \qquad
  \E[N\mid\bm\beta,\text{data}]=\sum_{i=1}^{n}\frac{1}{p_i(\bm\beta)},
\end{equation}
and by part (ii) of the lemma the marginal posterior of $\bm\beta$ is proportional to $p(\bm\beta)L_c(\bm\beta)$. The Horvitz--Thompson estimator is therefore the conditional posterior mean of $N$ given $\bm\beta$. Posterior draws of $N$ combine uncertainty about $\bm\beta$ with geometric ``unseen twins'' of each recorded incident. Two caveats carry over from the frequentist analysis. First, profiles that no list ever records are not represented. The posterior refers to the incidents whose profiles lie in the observed support, and nothing in \eqref{eq:models-ht} extrapolates beyond it. Second, a recorded incident with small $\hat p_i$ carries a large weight. Such incidents dominate both the estimate and its variance and should be inspected individually; a common cause is a linkage error that has split one incident into several singletons.

\paragraph{Random effects} Lists differ across places and over time in ways that $\bm x_i$ does not capture. A group-level random effect accommodates this:
\begin{equation}\label{eq:models-re}
  \logit p_{ik}=\beta_{0k}+\bm x_i^{\top}\bm\beta_k+b_{g(i),k},\qquad \bm b_g=(b_{g1},\dots,b_{gK})^{\top}\overset{\text{iid}}{\sim}N_K(\bm0,\bm\Sigma_b),
\end{equation}
where $g(i)$ is the state, or the state--year, of incident $i$. Given $\bm b$, incidents remain independent and $L_c$ keeps its form, so $\bm b$ is integrated out by MCMC. The correlation matrix of $\bm\Sigma_b$ has a direct reading. A positive correlation between the police and NGO effects means that states whose police record more also have NGOs that record more. This is dependence between domains, which a pooled model would misread as dependence between lists. An incident-level random effect, shared by all lists, gives the logistic-normal heterogeneity model \citep{coull1999}. Its conditional likelihood requires an integral over the random effect in both numerator and denominator. Its estimate of $N$ also depends heavily on the assumed shape of the random-effect distribution in the left tail, where the data are silent \citep{link2003}. We therefore model unobserved incident-level heterogeneity with latent classes and observed heterogeneity with covariates. Categorical covariates, such as offence type, target and setting, enter the latent-class model \eqref{eq:lcmcr} by stratification. Either the model is fitted separately within each stratum, or the class profiles $\lambda_{j\cdot}$ are shared across strata while the class weights vary by stratum under a hierarchical prior. \Cref{lem:models-negbin} then applies stratum by stratum. \Cref{sec:models-sae} combines latent classes and covariates with spatial and temporal pooling.

\subsection{A spatiotemporal small-area model}\label{sec:models-sae}

Antisemitic incidents are not spread evenly across the country. In many state--years the lists together record only a handful of incidents, and many cells of the capture table are empty. A separate MSE fit in each state--year would be unstable, and a single national fit would hide the variation across places and years that matters for policy. We therefore model all state--years jointly, so that sparse domains borrow strength from their neighbours in space and time. Let $s=1,\dots,S$ index states and $t=1,\dots,T$ index years, so that the domain is $d=(s,t)$. Let $P_{st}$ be an exposure, such as the resident population; the choice of exposure is discussed in \cref{sec:plan}. The model is
\begin{equation}\label{eq:smallarea}
\begin{aligned}
  n_{st,\h}\mid\theta_{st},\bm\pi_{st}
    &\overset{\text{ind}}{\sim}\Pois\big(P_{st}\,\theta_{st}\,\pi_{st,\h}\big),\quad \h\in\Hset,
  \qquad N_{st}=\textstyle\sum_{\h\in\Hset}n_{st,\h},\\
  \log\theta_{st}&=\mu+\bm x_{st}^{\top}\bm\beta+b_s+\gamma_t+\delta_{st},\\
  b_s&=\sigma_b\big(\sqrt{\phi}\,u_s+\sqrt{1-\phi}\,v_s\big),
  \qquad \bm u\sim\mathrm{ICAR}_\star(\mathcal G),\quad v_s\overset{\text{iid}}{\sim}N(0,1),\\
  \gamma_t-2\gamma_{t-1}+\gamma_{t-2}&\overset{\text{iid}}{\sim}N(0,\sigma_\gamma^2),
  \qquad \delta_{st}-\delta_{s,t-1}\overset{\text{iid}}{\sim}N(0,\sigma_\delta^2),\\
  \pi_{st,\h}&=\sum_{j=1}^{J}\pi_j\prod_{k=1}^{K}p_{st,jk}^{h_k}(1-p_{st,jk})^{1-h_k},
\end{aligned}
\end{equation}
with the capture probabilities
\begin{equation}\label{eq:models-capture}
  p_{st,jk}=c_{st,k}\,\operatorname{expit}\big(\alpha_{jk}+\eta_{ks}+\xi_{kt}+\zeta\,\ind{k=1,\ t\ge T^\star_s}+\bm x_{st}^{\top}\bm\beta_k\big),
\end{equation}
where $\bm\eta_s=(\eta_{1s},\dots,\eta_{Ks})^{\top}\overset{\text{iid}}{\sim}N_K(\bm0,\bm\Sigma_\eta)$ and $\xi_{k\cdot}$ is an independent first-order random walk for each list. By \cref{lem:fw-factor}(a), the first line of \eqref{eq:smallarea} is equivalent to $N_{st}\sim\Pois(P_{st}\theta_{st})$ together with a multinomial split of the $N_{st}$ incidents across histories with the latent-class probabilities of the last line. The incidence model is a disease-mapping model. Its spatial term is the BYM2 reparameterisation \citep{riebler2016} of the convolution prior of \citet{besag1991}. Here $\mathcal G$ is the graph of states that share a border, $\mathrm{ICAR}_\star$ is the intrinsic conditional autoregression scaled so that the geometric mean of its marginal variances is one \citep{sorbye2014}, $\sigma_b$ is the total standard deviation of the state effects and $\phi\in[0,1]$ is the share of that variance that is spatially structured. The national trend $\gamma_t$ is a second-order random walk. The interaction $\delta_{st}$ lets each state deviate from the national trend along its own random walk; this is the type~II interaction of \citet{knorrheld2000}, and types~I and~IV serve as sensitivity analyses. The latent classes, with weights $\pi_j$ and baseline profiles $\alpha_{j\cdot}$, are shared by all domains. Domain-specific capture shifts $\eta_{ks}$ and $\xi_{kt}$ move every class's capture probability for list $k$ together on the logit scale. We use a finite $J$ with a sparse Dirichlet prior on $(\pi_j)$, which empties superfluous classes \citep{rousseau2011}, because the model is fitted by Hamiltonian Monte Carlo.

\paragraph{Police coverage as data, and the NIBRS break} The police list can record an incident only if the incident falls in the jurisdiction of an agency that submits hate-crime data for the year. Let $c_{st,1}\in[0,1]$ be the share of state $s$'s population served by such agencies in year $t$. This share is computable from the FBI's agency participation data (\cref{sec:sources}), and we enter it in \eqref{eq:models-capture} as a known multiplicative offset. The offset is exact under two assumptions. First, an incident in a non-participating jurisdiction cannot enter list~1. Second, within a state--year, incidents are spread across covered and uncovered jurisdictions in proportion to population. The second assumption can be relaxed. One option replaces $c_{st,1}$ by an incidence-weighted coverage $c_{st,1}e^{\epsilon_{st}}$ (truncated at one) with $\epsilon_{st}\sim N(0,\tau_c^2)$; another computes coverage separately for urban and rural agencies.

The offset matters most around the national program's move from summary to incident-based (NIBRS) collection \factcheck{the FBI stopped accepting summary (SRS) submissions beginning with the 2021 data year}. After that move, agency participation fell sharply in some states for some years \factcheck{notably the 2021 data year}. This is a structural break in list-1 coverage. Without the offset, a model whose police capture probabilities evolve smoothly cannot represent the break. It must attribute the fall in police counts partly to a fall in $N_{st}$ and partly to a fall in police recording, in proportions set by the prior. With $c_{st,1}$ in the model, the break is absorbed by data, and both the remaining police capture effects and the incidence trend $\gamma_t$ remain smooth. The offset does not capture changes in practice among agencies that \emph{did} participate, for example differences in how the two reporting systems record bias motivation. The term $\zeta\,\ind{t\ge T^\star_s}$, where $T^\star_s$ is the first year in which state $s$ reported through NIBRS only, absorbs a level shift of that kind. Since the NGO and community lists are not affected by the transition, the break parameter $\zeta$ is identified from the change in police capture of the incidents those lists record.

\paragraph{The other lists: hierarchical capture} For $k\ge2$, $c_{st,k}\in\{0,1\}$ records whether list $k$ operates in state $s$ in year $t$. A list that does not operate somewhere has a structural zero there. This replaces \cref{asm:fw-scope}(b), which requires every list to be able to record every incident, by a modelled footprint. The state effects $\bm\eta_s$ share the covariance $\bm\Sigma_\eta$ across states. This pools each list's state-specific capture rate toward its national level, and the off-diagonal elements describe the tendency of police and community organisations to be active, or inactive, in the same states. A state with few incidents thus has its capture rates, and hence its dark-figure ratio $D_1$, shrunk toward the national pattern unless its own capture table says otherwise.

\paragraph{Borrowing strength} Three mechanisms pool information across domains. First, for incidence, the ICAR component has full conditional prior $u_s\mid\bm u_{-s}\sim N(\bar u_{\partial s},1/d_s)$, where $\bar u_{\partial s}$ is the mean of the effects of the $d_s$ neighbours of state $s$ (before scaling). To a Gaussian approximation, if the capture table of state $s$ yields a local log-incidence estimate with sampling variance $v_s$, the structured effect is a precision-weighted compromise between that estimate and the neighbourhood mean, with weights $v_s^{-1}$ and a prior precision proportional to $d_s/(\sigma_b^2\phi)$. When $v_s$ is large, as it is for a sparse state, the neighbours dominate. The temporal terms do the same across adjacent years. Second, for capture, the hierarchical effects $\bm\eta_s$ shrink sparse states toward the national capture profile. Third, the latent reportability types and their capture profiles are learned from the whole country. A sparse state--year therefore needs only to determine its incidence and its capture shifts, not a full latent structure.

By Poisson thinning, given the parameters the unobserved count $n_{st,\zero}\sim\Pois(P_{st}\theta_{st}\pi_{st,\zero})$ is independent of the observed cells. Posterior draws of $N_{st}=n_{st}+n_{st,\zero}$ and of any aggregate, such as the national total $\sum_sN_{st}$ or a regional trend, are obtained by drawing $n_{st,\zero}$ at each iteration. These draws carry the full posterior correlation across domains. Borrowing strength reduces variance, but it does not identify $\omega$ in any domain. It transfers the identifying restriction of the latent-class model from the nation to each state. A state whose reporting environment is atypical in a way the covariates and random effects miss, for instance one where the NGO list is fed by police referrals (\cref{prop:referral}), will be shrunk toward a national pattern that does not apply to it. Domain-level sensitivity analysis with \cref{cor:bounds} therefore remains part of the analysis plan.

\subsection{Propagating linkage uncertainty}\label{sec:models-link}

The capture tables above are functions of a linkage structure. Let $R$ denote all records with their attributes, let $\Link$ be a partition of $R$ into incidents (\cref{sec:linkage}), and let $\bm n(\Link)$ be the collection of domain capture tables that $\Link$ induces. A fully Bayesian treatment specifies a single joint model $p(R\mid\Link,\psi_L)\,p(\Link\mid N,\psi)\,p(N,\psi)\,p(\psi_L)$ for records, linkage, population and capture, as in \citet{tancredi2011} and \citet{tancredi2020}. Here $\psi_L$ are the linkage parameters and $\psi$ the capture parameters. In this joint model the population model informs which records are linked. We use instead the modular (``cut'') posterior \citep{bayarri2009,plummer2015}, which \citet{sadinle2018} used to propagate linkage uncertainty into population-size estimates of human rights violations:
\begin{equation}\label{eq:models-cut}
  p_{\mathrm{cut}}(N,\psi\mid R)=\int p\big(N,\psi\mid\bm n(\Link)\big)\,p_L(\Link\mid R)\,d\Link,
\end{equation}
where $p_L(\Link\mid R)$ is the posterior of a linkage model fitted without reference to the capture model. The cut has three practical advantages. Each module can be checked separately. Misspecification of the capture model cannot distort the linkage. And one set of linkage draws serves every MSE model in \cref{sec:models-lcmcr,sec:models-cov,sec:models-sae}, as well as the estimators of \cref{sec:identification}. The cost is that the weak information the population model carries about which particular pairs match is discarded. Under \eqref{eq:models-cut} the posterior mean is $\E_L\{\E[N\mid\bm n(\Link)]\}$, and the posterior variance splits into a within-linkage and a between-linkage component, as in \cref{prop:totalvar}. The plug-in posterior $p(N\mid\bm n(\Linkhat))$ drops the second component, and it is centred at a biased value whenever $\Linkhat$ systematically splits or merges incidents (\cref{thm:linkbias}). When only a point linkage $\Linkhat$ with estimated error rates $(\hat\delta,\hat\varepsilon)$ is available, the corrected table of \cref{cor:linkcorrect} is a fallback. Agreement between the cut posterior and the corrected plug-in estimate is a useful diagnostic of the linkage-error model. \Cref{alg:pipeline} gives the end-to-end procedure, and \cref{fig:pipeline} shows it schematically.

\begin{algorithm}[t]
\caption{End-to-end estimation of the dark figure with propagated linkage uncertainty.}\label{alg:pipeline}
\begin{algorithmic}[1]
\Require Lists $R_1,\dots,R_K$ with harmonised attributes; blocking rule $\mathcal B$; numbers of linkage draws $M$ and MSE draws $L$; an MSE model from \cref{sec:models-lcmcr,sec:models-cov,sec:models-sae}.
\Ensure Draws $\{N_d^{(m,\ell)}\}$ from the cut posterior \eqref{eq:models-cut} for every domain $d$.
\State \textbf{Deduplicate} within each list $R_k$, so that each incident contributes at most one record per list.
\State \textbf{Block}: form candidate pairs $\{(r,r')\in R_k\times R_{k'}:k<k',\ \mathcal B(r,r')=1\}$ and their comparison vectors $\gamma_{rr'}$.
\State \textbf{FS-EM}: estimate $m(\cdot)$ and $u(\cdot)$ by \cref{alg:fsem}; compute the weights $w_{rr'}=\log\{m(\gamma_{rr'})/u(\gamma_{rr'})\}$.
\State (Optional) \textbf{Clerical review} of the most uncertain pairs by \cref{alg:review}; fix the reviewed pairs.
\For{$m=1,\dots,M$}
  \State Draw $\Link^{(m)}\sim p_L(\Link\mid R)$ by MCMC over partitions with at most one record per list per incident.
  \State Form the capture tables $\bm n_d^{(m)}=\bm n_d(\Link^{(m)})$ for every domain $d$.
  \State Run the MSE sampler on $\{\bm n_d^{(m)}\}_d$; keep $L$ draws $N_d^{(m,\ell)}=n_d^{(m)}+n_{d,\zero}^{(m,\ell)}$, $\ell=1,\dots,L$.
\EndFor
\State \textbf{Mix}: pool the $ML$ draws; report posterior summaries and the between- and within-linkage variance components of \cref{prop:totalvar}.
\end{algorithmic}
\end{algorithm}

Three remarks on \cref{alg:pipeline}. First, line~6 needs a sampler over $K$-partite partitions, not just bipartite matchings. The Fellegi--Sunter likelihood ratios $w_{rr'}$ enter its target, and the constraint of at most one record per list per incident (after deduplication) is enforced in its proposals; \citet{sadinle2017} and \citet{steorts2016} give samplers of this kind. Plugging in the FS-EM estimates of $m$ and $u$ is an empirical-Bayes shortcut that slightly understates linkage uncertainty. A fully Bayesian linkage model samples them as well. Second, the draws $\Link^{(m)}$ come from a Markov chain, so the between-linkage component is estimated from $M_{\mathrm{eff}}<M$ effectively independent tables. \Cref{app:computation} gives the Monte Carlo error of the pooled estimate and the choice of $L$ that minimises it for a fixed budget. Third, the MSE step in line~8 is embarrassingly parallel across $m$. Its output can be reused for the funnel decomposition of \cref{sec:funnel}, which needs the linkage-specific overlap between police and NGO records.
